%% halaman-judul.tex
%% Halaman Judul (Sampul Dalam)
%% Format sama dengan cover.tex tapi tanpa background kuning
%% Font: Sans-serif (TeX Gyre Heros) seperti cover
%%
%% EDIT LANGSUNG TEKS DI BAWAH INI SESUAI DATA TESIS ANDA
\newpage
\thispagestyle{plain}
% Temporarily switch to sans-serif font for this page
{
\sffamily % Switch to sans-serif font

% All layout using TikZ for precise positioning
\begin{tikzpicture}[remember picture,overlay]
    % ITS Logo - positioned at top left corner
    \node[anchor=north west] at ([xshift=1.5cm,yshift=-2cm]current page.north west) 
        {\includegraphics[width=3.5cm]{figures/logo-its.png}};
    
    % Blue stripe - full width
    \fill[ITSBlue] ([yshift=-6cm]current page.north west) rectangle ([yshift=-6.5cm]current page.north east);
    
    % Content - positioned below blue stripe
    % UBAH NILAI yshift=-8cm UNTUK MENGATUR JARAK DARI ATAS
    \node[anchor=north west, text width=15cm] at ([xshift=2.0cm,yshift=-8cm]current page.north west) {
        \begin{minipage}{15cm}
        \begin{onehalfspace}
        
        %% ===== KODE TESIS =====
        {\fontsize{12pt}{14pt}\selectfont\bfseries Tesis-Sidang Akhir - EF235401}
        
        \vspace{0.4cm}
        
        %% ===== JUDUL TESIS (18pt bold) =====
        %% Edit judul di sini
        {\fontsize{18pt}{22pt}\selectfont\bfseries 
        Judul Tesis (Ditulis \textit{Capitalize Each Word} – sesuai EYD)}
        
        \vspace{1.8cm}
        
        %% ===== NAMA DAN NRP MAHASISWA =====
        {\fontsize{12pt}{14pt}\selectfont\bfseries 
        NAMA MAHASISWA\\
        NRP MAHASISWA}
        
        \vspace{1cm}
        
        %% ===== DOSEN PEMBIMBING =====
        {\fontsize{12pt}{14pt}\selectfont\bfseries DOSEN PEMBIMBING}
        
        \vspace{0.15cm}
        
        {\fontsize{12pt}{14pt}\selectfont\bfseries 
        Nama Dosbing 1\\
        Nama Dosbing 2}
        
        \vspace{1cm}
        
        %% ===== INFORMASI PROGRAM =====
        {\fontsize{12pt}{14pt}\selectfont\bfseries 
        PROGRAM MAGISTER\\
        BIDANG KEAHLIAN .................(DIISI RMK)\\
        PROGRAM STUDI S2 TEKNIK INFORMATIKA\\
        DEPARTEMEN TEKNIK INFORMATIKA\\
        FAKULTAS TEKNOLOGI ELEKTRO DAN INFORMATIKA CERDAS\\
        INSTITUT TEKNOLOGI SEPULUH NOPEMBER\\
        SURABAYA\\
        TAHUN}
        
        \end{onehalfspace}
        \end{minipage}
    };
\end{tikzpicture}

} % End of sffamily